\section{Reading Comprehension Questions}

\noindent
{\em From Section~\ref{sec:graphTypes}}
\begin{requ}
Draw an example of each of the following:
\begin{lenum}
\item An weighted undirected pseudograph
\item An unweighted directed multigraph
\item A network
\end{lenum}
\begin{requanswer}
(a) answers will vary. Your graph should have numbers (weights) on every edge, 
no arrows on the edges, and can contain loops (edges from a vertex to itself).
(b) answers will vary. You graph should not have any numbers on the edges,
the edges should have arrows, and there may be repeated edges--that is, there
might be two different edges from some vertex $u$ to another vertex $v$.
(c) answers will vary. A network is just a weighted directed graph. 
So your graph should have numbers on the edges and every edge should have an arrow.
It should not have loops or repeated edges.
\end{requanswer}
\end{requ}

\begin{requ}
Give an example of a problem that might be modeled using the following types of graphs.
Make sure it is clear what the vertices and edges represent.
\begin{lenum}
\item A network
\item An directed weighted multigraph. 
\item An unweighted undirected pseudograph
\end{lenum}
\begin{requanswer}
(a) A road system, where the vertices are intersections and the edges are the segments of road between intersections.
The weights on the edges might be distances, speed limits, expected time to traverse, etc.
(b) A ski trail map, where the vertices are intersections and the edges are the segments of trail between intersections.
The edges are directed because on som trails you are only allowed to go one direction.
Since sometimes a trail splits and comes back together, multiple edges are allowed between vertices.
The weights can be distances or difficulty ratings.
(c) Representing connections on a social media app, where it is assumed that being connected is two-way
(e.g. on Facebook when you are friends versus on Instagram where following is one-way), and where you are
allows to connect to yourself (I actually do not know of a social media site that allows this, but I suppose it is possible).
\end{requanswer}
\end{requ}

\begin{requ}
Prove that a graph with $n$ vertices has at most $\binom{n}{2}$ edges.
(There are several possible ways to prove this. You can use counting techniques or induction, for instance.)
\begin{requanswer}
The easiest proof is to realize that edges are just pairs of vertices. There are $n$ vertices. How many ways are there of
choosing pairs of vertices? You are choosing 2 things out of $n$ things, so $\binom{n}{2}$.

Here is a proof by induction: A graph with 2 vertices has $1=\binom{2}{2}$ possible edge, so it holds for the base case
(We could use 1 vertex as a base case, but it is more confusing so we start at 2).
Assume a graph with $k-1$ vertices, where $k\geq 3$ has $\binom{k-1}{2}$ possible edges. 
If you add a vertex, it can be connected to each of
the $k-1$ vertices, so a graph with $k$ vertices has 
$\binom{k-1}{2}+(k-1) = \frac{(k-1)(k-2)}{2}+(k-1) = (k-1)\left(\frac{k-2}{2} + 1\right) 
= (k-1)\left(\frac{k-2+2}{2}\right)
=\frac{(k-1)(k)}{2}=\binom{k}{2}$ possible edges.
Since the formula is true for $k=2$, and whenever it is true for $k-1$ it is true for $k$, it is true for all $n\geq 2$ by induction.
\end{requanswer}
\end{requ}

\noindent
{\em From Section~\ref{sec:graphTerm}}
\begin{requ}
Draw an unconnected graph such that one component contains a cycle of length 4 and another component is a tree.
\begin{requanswer}
answers will vary, but here are two examples:  $\square \, |$ \, or \,  $\boxtimes \, \wedge$.
\end{requanswer}
\end{requ}

\begin{requ}
\begin{lenum}
\item
How many edges does a tree with $n\geq 2$ vertices have? 
Draw a few trees of various sizes and you should see an obvious pattern.
\item
(a bit challenging) Prove that your formula is correct. 
(Hint: Use induction. But you have to be a little careful in how you do it.
Also, you may assume that every tree contains at least one vertex with degree 1.)
\end{lenum}
\begin{requanswer}
(a) $n-1$. 
(b) Clearly a tree with $2$ vertices has $2-1=1$ edges.
Assume all trees with $n-1$ vertices have $n-2$ edges.
Let $T$ be a tree with $n>2$ vertices. Since it is a tree, it contains at least one vertex $v$ of degree 1.
Let $T^\prime$ be the tree $T$ with vertex $v$ deleted.
Then $T^\prime$ has $n-1$ vertices and 
is clearly still a tree since removing a vertex of degree 1 cannot either add a cycle or disconnect the graph.
Thus $T^{\prime}$ has $n-2$ edges. But it was created from $T$ by removing a single vertex and edge.
Therefore $T$ has $n-1$ edges.
Since the formula holds for $n=2$ and whenever it holds for $n-1$ it holds for $n$, 
every tree with $n\geq 2$ vertices has $n-1$ edges.
\end{requanswer}
\end{requ}

\begin{requ}
answer the following questions about graph $L$ below.

\newcounter{mygraphquestioncounter}

\noindent
\begin{minipage}[b]{.5\textwidth}
\begin{lenum}
\item Is $L$ weighted or unweighted? 
\item Is $L$ directed or undirected? 
\item Are $e$ and $x$ adjacent? Are $f$ and $b$ adjacent?
\item Is $L$ connected?
\item How many vertices does $L$ have?
\item How many edges does $L$ have?
\item What is $\deg(v)$? $\deg(c)$? 
\setcounter{mygraphquestioncounter}{\value{enumi}}
\end{lenum}
\end{minipage}
\hfill
%\begin{minipage}{.45\textwidth}
\vspace*{0in}
\sputfig{graphs/fig/lemke_max}{.75}
%\end{minipage}

\begin{lenum}
\setcounter{enumi}{\value{mygraphquestioncounter}}
\item What is $\sum_{v\in L} \deg(v)$? 
Does this number seem to be related to your answer from (e)? Explain.
\item Draw a spanning tree of $L$. How many edges does it have?
Does your spanning tree contain any cycles? Explain why or why not.
\item What is the minimum number of edges that you can remove to make the graph disconnected
(that is, not connected)? Which ones?
\item Find a cycle of length 3 in $L$. Then find one of length 4. Repeat for 5, 6, 7, and 8.
\end{lenum}
\begin{requanswer}
(a) unweighted.
(b) undirected.
(c) $e$ and $x$ are not adjacent, but  $f$ and $b$ are adjacent. 
(d) $L$ is connected.
(e) 8.
(f) 17.
(g) $\deg(v)=2$ and $\deg(c)=5$. 
(h) 34. It is twice the number of edges which makes sense because of the Handshaking Lemma.
(i) answers will vary, but if you drew a subgraph of $L$ that contains exactly 7 edges and contains no cycles,
then it is a spanning tree. Of course it cannot contain cycles because it is a tree.
(j) If you remove $(e,v)$ and $(f,v)$, then $v$ is disconnected from the rest of the graph. No single edge
will disconnect the graph so 2 edges is the minimum number.
(k) answers will vary, but mine are $efv$, $bcdf$, $abcdx$, $axdfve$, $aefbcdx$, and finally $aevfbcdx$.
Notice in all of these, the vertices listed next to each other are adjacent and the first and last one on the list are adjacent.
\end{requanswer}
\end{requ}

\noindent
{\em From Section~\ref{sec:specialGraphs}}
\begin{requ}
Draw $K_6$.
\begin{requanswer}
Here are the graphs for the next 3 questions.
The vertices labels are optional.

\sputfig{graphs/fig/K6}{.8}
\end{requanswer}
\end{requ}
\begin{requ}
Draw $C_8$.
\begin{requanswer}
See above.
\end{requanswer}
\end{requ}
\begin{requ}
Draw $P_5$.
\begin{requanswer}
See above.
\end{requanswer}
\end{requ}

\begin{requ}
Draw $Q_5$. Just kidding. That would be a bit difficult to visualize. Instead, describe how you could
construct $Q_5$ recursively. For instance, can you see how to go from $Q_1$ to $Q_2$?
And from $Q_2$ to $Q_3$? 
And from $Q_3$ to $Q_4$?  
Once you observe the pattern it is pretty straightforward to see how
to construct $Q_{k+1}$ from $Q_k$.
\begin{requanswer}
Notice to draw $Q_2$, you can draw two copies of $Q_1$ (i.e. two lines) and connect corresponding
vertices (i.e. if you draw the two lines vertically, connect the top vertices and the bottom vertices) to
make a square (notice that $Q_2$ is the same as $C_4$).
To draw $Q_3$ (a cube), you can draw two copies of $Q_2$ (squares) and connect corresponding vertices.
Using the same idea, to draw $Q_5$, 
I would draw 2 copies of $Q_4$, and then connect corresponding vertices between the two copies.
\end{requanswer}
\end{requ}

\begin{requ}
Give a partition of the vertices of $Q_3$ to show that it is bipartite.
In other words, which vertices go in $V_1$ and which go in $V_2$?
Use 3-bit numbers to list the vertices (since that is the natural way to construct the graph).
\begin{requanswer}
Notice that two vertices are adjacent iff they differ by exactly 1 bit.
Also notice that if two numbers have the same parity, they cannot differ by exactly 1 bit.
So we can pick $V_1=\{000,011,101,110\}$, and $V_2=\{001,010,101,111\}$.
The numbers in $V_1$ have even parity, and the numbers in $V_2$ have odd parity.
So within each subset, none of the numbers differ by exactly one bit since all of the numbers in each
set have the same parity. Thus, this is a valid partition.
\end{requanswer}
\end{requ}

\begin{requ}
Draw $K_{3,5}$.  Then draw a graph $G$ such that $G$ is a subgraph of $K_{3,5}$.
\begin{requanswer}
Here is $K_{3,5}$ and one possible subgraph.

\sputfig{graphs/fig/K3_5}{1}
\end{requanswer}
\end{requ}

\noindent
{\em From Section~\ref{sec:handshaking}}
\begin{requ}
Give an informal proof of Theorem~\ref{thm:handshake_lemma}.
That is, argue why it makes sense by talking about edges, degrees, and vertices.
\begin{requanswer}
Every edge has two endpoints. 
Each endpoint adds one to the degree of the vertex it is incident with.
So if you sum the degree of all of the vertices, it should be twice the number of edges.
\end{requanswer}
\end{requ}

\begin{requ}
You are at a party with some friends and one of them claims 
``I just did a quick count, and it turns out that at this party,
there are an odd number of people who have shaken hands with an odd number of other people at the party.''
Prove or disprove that this friend is correct.
\begin{requanswer}
This is a simple application of the Handshaking Lemma.
We create a graph as follows: 
People are vertices, and there is an edge between two people if they have shaken hands.
So the statement would imply that in the graph there are an odd number of vertices of odd degree.
This contradicts Corollary~\ref{cor:hand} (which is a corollary of the Handshaking Lemma).
Thus, the friend is incorrect.
\end{requanswer}
\end{requ}

\noindent
{\em From Section~\ref{sec:graphRep}}
\begin{requ}
\begin{lenum}
\item If a graph has very few edges, which representation is a better choice if space is the only consideration? Explain.
\item If a graph has many edges, which representation is a better choice if space is the only consideration? Explain.
\end{lenum}
\begin{requanswer}
(a) Recall that an adjacency list requires on the order of $\Theta(n+m)$ memory, whereas for the
adjacency matrix it would be $\Theta(n^2)$. 
If the graph has few edges, then $m$ is much smaller than $n^2$, so $\Theta(n+m)$ would be smaller
and the adjacency list would be appropriate.
(b) If the number of edges is larger, then either might be appropriate, depending on how large. If 
$m\approx n^2$, then we are comparing $\Theta(n+m)=\Theta(n+n^2)=\Theta(n^2)$ with $\Theta(n^2)$,
so there is minimal difference between the two. But if $m$ is large but smaller than $n^2$, the adjacency list
might be the better choice.
\end{requanswer}
\end{requ}

\begin{requ}
Are space considerations actually that important? In other words, practically speaking, if you can store a graph 
using one of the representation, can you store it in the other without worrying too much about space? Explain.
(This is an important question, so think carefully about it!) (Hint: Think about storing the graph of friends on Facebook
or another social media site.)
\begin{requanswer}
If you are storing small graph (e.g. hundreds of vertices or less), it probably does not matter a whole lot.
But imagine storing the Facebook friendship graph. As of 2021, there are 2.85 billion Facebook users
and each has an average of 350 friends (as of 2019 it was about 338, so this number should be close).
Since we are talking about exact numbers, we will compare without using $\Theta$ notation.
Recall that an adjacency list  takes $\Theta(n+m)$ space and an adjacency matrix takes $\Theta(n^2)$ time.
We will just treat these as $n+m$ and $n^2$.
In our example, $n=2,850,000,000$, and $m=2,850,000,000*350=997,500,000,000$.
So an adjacency list would take about $n+m=2,850,000,000+997,500,000,000=1,000,350,000,000$ space.
An adjacency matrix would take about $n^2=2,850,000,000^2=8,122,500,000,000,000,000$ space.
In case it isn't clear, the adjacency matrix takes about 8,119,658 times as much space!
To be more precise, assuming it takes 32 bits to store each number in an adjacency list or matrix 
(and that is pushing it), 
the adjacency list requires about 4 TB (terabtyes) of space, which is doable if you want to fill up the
majority of the hard drive on a very new top-of-the-line computer.
On the other hand, the adjacency matrix requires about 32.49 EB (exabytes).
The largest hard drive you can currently buy is about 18 TB, so you would need about 1,804,369
hard drives to store the adjacency matrix. (Even if you encoded the matrix densely, using only 1 bit per entry,
it would still require about 56,387 hard drives.)
So yes, it does matter.
\end{requanswer}
\end{requ}

\begin{requ}
Given an {\em adjacency list} representation of a graph with $n$ vertices and $m$ edges, 
how long do the following operations take?
\begin{lenum}
\item Determine whether or not $(u,v)\in E$.
\item Determine $\deg(u)$.
\item Iterate over the neighbors of vertex $u$ (assume $u$ has $k$ neighbors). 
\end{lenum}
\begin{requanswer}
(a) Arbitrarily pick either $u$ or $v$ and check its list for the other--we will look at $u$'s list.
Since it might have to traverse the entire list of the neighbors of $u$, it would be $O(\deg(u))$,
which might be as large as $n-1$. So $O(n)$ in the worst case (although $O(\deg(u))$ is more precise).
I use big-$O$ notation on this one because it is possible it finds it sooner.
(b) Either $\Theta(\deg(u))$ if it has to traverse the list and count, 
or $\Theta(1)$ if this is maintained in the data structure.
(c) $\Theta(\deg(u))=\Theta(k)$.
\end{requanswer}
\end{requ}

\begin{requ}
Given an {\em adjacency matrix} representation of a graph with $n$ vertices and $m$ edges, 
how long do the following operations take?
\begin{lenum}
\item Determine whether or not $(u,v)\in E$.
\item Determine $\deg(u)$.
\item Iterate over the neighbors of vertex $u$ (assume $u$ has $k$ neighbors). 
\end{lenum}
\begin{requanswer}
(a) $\Theta(1)$ since it can just look at the $(u,v)$ entry of the matrix.
(b) $\Theta(n)$ since it has to look through an entire row of the matrix to determine which vertices are neighbors.
(c) Same answer and reason as (b).
\end{requanswer}
\end{requ}

\begin{requ}
\begin{lenum}
\item If adding and removing {\em edges} is an important operation, is one of the representations a better choice? Explain.
\item If adding and removing {\em vertices} is an important operation, is one of the representations a better choice? Explain.
\end{lenum}
\begin{requanswer}
(a) Either one works fine, but it is more efficient with an adjacency matrix since an edge $(u,v)$ can be added and removed in 
constant time by just changing the matrix entry $(u,v)$ to a 0 or 1. 
For an adjacency list, to remove $(u,v)$, you would have to find $u$ on $v$'s list and $v$ on $u$'s list and then remove them from the lists, so it would take longer.
(b) Adjacency list by far. If you add a vertex, you can just add an adjacency list for it to your current list. With a matrix,
you need to create an entirely new matrix with one more row and column and copy all of the entries, so it is not very efficient.
Similar problems exist when removing vertices.
\end{requanswer}
\end{requ}


\noindent
{\em From Section~\ref{sec:trails}}

\begin{requ}
\begin{lenum}
\item Is $K_5$ Eulerian? Explain.
\item Is $K_6$ Eulerian? Explain.
\item Is $Q_3$ Eulerian? If so, number the edges of $Q_3$ in order to demonstrate the Euler tour. If not, explain why not.
\item Is $Q_4$ Eulerian? If so, number the edges of $Q_3$ in order to demonstrate the Euler tour. If not, explain why not.
\item Is $Q_3$ Hamiltonian? If so, draw a Hamiltonian cycle on $Q_3$. If not, explain why not.
\item Is $Q_4$ Hamiltonian? If so, draw a Hamiltonian cycle on $Q_4$. If not, explain why not.
\item Is asking if $C_7$ is Hamiltonian a stupid question? Explain.
\end{lenum}
\begin{requanswer}
(a) Since every vertex of $K_5$ has degree 4, it is Eulerian by Theorem~\ref{thm:Eulerian}.
(b) Since every vertex of $K_6$ has degree 5, it is {\em not} Eulerian by Theorem~\ref{thm:Eulerian}.
(c) Since every vertex of $Q_3$ has degree 3, it is {\em not} Eulerian by Theorem~\ref{thm:Eulerian}.
(d) Since every vertex of $Q_4$ has degree 4, it is Eulerian by Theorem~\ref{thm:Eulerian}.
Here is one of many possible orderings of the edges that form an Eulerian tour:

\sputfig{graphs/fig/hypercube4_Euler}{1}

\noindent
(e) and (f) are both are Hamiltonian. Here is one possible Hamiltonian cycle for each:
%000, 100, 110, 010, 011, 111, 101, 001, 000
%0000, 0100, 0110, 0010, 0011, 0001, 0101, 0111, 1111, 1011, 1010, 1110, 1100, 1101, 1001,1000, 0000.

\sputfig{graphs/fig/hypercube3_Ham}{1}
\sputfig{graphs/fig/hypercube4_Ham}{.8}

\noindent
(g) Yes. Since it is just a cycle with all of the vertices, it is clearly a Hamiltonian cycle. 
\end{requanswer}
\end{requ}

\noindent
{\em From Section~\ref{sec:planar}}

\begin{requ}
Give a planar embedding of $K_{2,3}$. 
\begin{requanswer}
Here is the most obvious way to draw it: 

\sputfig{graphs/fig/K2_3}{1}
\end{requanswer}
\end{requ}

% Dang it. This is addressed at the end of the chapter.
%\begin{requ}
%Is each of the following graphs planar or not?
%If it is planar, give a planar embedding. Otherwise try to explain why you do not think it is possible.
%(Note: If it is planar, your planar embedding is a proof of that fact. 
%However, proving a graph is not planar is 
%not so easy. Thus, I am not asking for a proof if they are not planar---just 
%experiment with each graph long enough until you are relatively sure that it is not planar.
%But realize that being relatively sure is not the same as having a proof!)
%\begin{lenum}
%\item $K_{3,3}$
%\item $K_5$
%\end{lenum}
%\end{requ}

\begin{requ}
Does Theorem~\ref{thm:kurat} imply that if a graph with $v\geq 3$ vertices has fewer
than $3v-6$ edges that it is planar? Explain, using an example if appropriate.
\begin{requanswer}
It is not saying that. No planar graphs have more edges than that, but not every graph with fewer edges is planar.
For instance, $K_{3,3}$ has $v=6$ and $e=9$. So $e=9<12=3v-6$, but as we saw earlier, $K_{3,3}$ is not planar. 
\end{requanswer}
\end{requ}

\begin{requ}
\begin{lenum}
\item Prove that $Q_3$ is planar.
\item Prove that $Q_4$ is not planar.
\end{lenum}
\begin{requanswer}
(a) I am too lazy to do another drawing, but draw it as a box inside a box with the corners connected to each other and it is clearly planar.
(b) Notice that $v=16$, $e=32$, and $C_3$ is not a subgraph of $Q_4$. Then $e=32>28=2v-4$, so by part (b) of
Theorem~\ref{thm:kurat}, $Q_4$ cannot be planar.
\end{requanswer}
\end{requ}


\noindent
{\em From Section~\ref{sec:mst}}

\begin{requ}
Can a graph have a minimum spanning tree with weight 34 
and another with weight 35? Explain.
\begin{requanswer}
No. By definition, the one with weight 35 would not be a minimum spanning tree.
\end{requanswer}
\end{requ}

\begin{requ}
An Internet provider is installing an optical network in a rural town.  
They need to wire every house. 
The map below indicates the locations of the houses (dots).
Each grid line represents 1 mile.
They can route the cables anywhere they wish, 
and only need one connection to each house.
Since optical cable is really expensive, 
the goal is to minimize the amount of cable required.
An example: 
If they want to route from {\bf I} to {\bf J}, they could route
south 3 miles, and east 1 mile for a total of 4 miles, 
but it would be better to route in 
a straight line from {\bf I} to {\bf J} for a cost of 
$\sqrt{3^2+1^2}=\sqrt{10}\approx 3.2$ miles.

\sputfig{graphs/fig/mstExa}{.6}
\begin{lenum}
\item
Explain how this problem can be modeled as a graph problem.
What are vertices? Edges? Is it directed? Weighted?
\item
Explain why solving this problem reduces to finding a minimum
spanning tree of the graph.
\item 
You will be asked to solve this problem in the next part.
Should you use an algorithm like Prim's or like Kruskal's? 
Explain. (Note: I say an algorithm``like'' because when
humans run algorithms on graphs drawn on paper, they
typically do not do them exactly as a computer would.)
\item
Compute a solution to this problem and give the total
length of wire necessary.
\end{lenum}
\begin{requanswer}
(a) The dots are vertices and we can think of it being a complete
graph where the weight of each edge is the distance between
the two dots.
(b) We want to route to every house, but we want to minimize
wire. Thus, a spanning tree makes sense because if we have any 
cycles, we are wasting wire. Not only that, we want a minimum 
spanning tree to use the least amount of wire.
(c) Prim's algorithm makes the most sense because it is a lot
easy to start somewhere and find the closest point to the
current partial tree. On the other hand, with Kruskal's algorithm,
we would have to sort a list if $\binom{n}{2}$ edges.
(d) There is a unique solution with a cost of about 50.3 as shown below.

\sputfig{graphs/fig/mstExaSol}{.7}
\end{requanswer}
\end{requ}

\begin{requ}
Consider the following graph $G$ with cut $S=\{A,B,C\}$
and $V-S=\{D,E\}$.

\sputfig{graphs/fig/wgExa2}{.4}

\begin{lenum}
\item List all of the cross edges.
\item List all of the light edges.
\item Which of the following edges crosses the cut:
$(A,B)$, $(A,D)$, $(D,E)$, $(B,D)$.
\item Find a minimum spanning tree of $G$ and give its weight.
\item Does $G$ have more than one minimum spanning tree? Explain.
\end{lenum}
\begin{requanswer}
(a) $(A,D)$, $(A,E)$, $(B,D)$, and $(C,D)$.
(b) $(A,E)$ is the only light edge.
(c) $(A,D)$ and $(B,D)$ cross the cut.
(d) The path with vertices $C, B, A, E, D$ is a minimum spanning tree
of weight $10$.
(e) The minimum spanning tree is unique because it uses the 4 
lightest edges of the graph so there can't be one with smaller weight.
\end{requanswer}
\end{requ}

\begin{requ}
Can you think of a necessary condition for a graph to 
have more than one minimum spanning tree?
Note: a {\em necessary condition} means a condition 
that has to be present for something to be true, but it being
present does not always guarantee that it is true.
\begin{requanswer}
A graph must have at least some edges of repeated weight
in order for there to be multiple minimum spanning trees.
\end{requanswer}
\end{requ}

